\chapter{Как узнать кошку}
% полная версия: chapters/12_recognition.tex

\pencilfig{12}{\foreach \cx/\cy/\s in {0/0/1,2.3/0.35/0.55} {
 \begin{scope}[shift={(\cx,\cy)}, scale=\s]
  \draw[pen, wash=pcAmber] (0,0) circle (0.45);
  \draw[pen] (-0.4,0.2) -- (-0.32,0.66) -- (-0.08,0.44); \draw[pen] (0.4,0.2) -- (0.32,0.66) -- (0.08,0.44);
  \fill[pcGreen] (-0.15,0.08) circle (0.05); \fill[pcGreen] (0.15,0.08) circle (0.05);
  \draw[pcGray, line width=0.4pt] (-0.1,-0.12) -- (-0.55,-0.05); \draw[pcGray, line width=0.4pt] (0.1,-0.12) -- (0.55,-0.05);
 \end{scope}}
\draw[pcGray, dashed, line width=0.5pt, ->] (0.55,0) -- (3.6,-0.35); \draw[pcGray, dashed, line width=0.5pt, ->] (2.65,0.35) -- (3.6,-0.2);
\node[lbl, anchor=west] at (3.65,-0.28) {кошка};}{Большая кошка и маленькая --- совсем разные картинки на датчиках глаза, а ответ один.}

Кошка на картинке может быть слева или справа, близко или далеко, на солнце или в тени. Каждый раз на датчики глаза падает совсем другой узор, а вы узнаёте её за доли секунды. В этом и состоит трудность распознавания: одно и то же нужно узнать в тысячах разных видов.

\section*{Конвейер}

Мозг решает задачу конвейером из примерно десятка ступеней, и сигнал проходит его за полторы десятых секунды. На каждой ступени делаются две вещи. Сначала --- «и»: нейрон срабатывает, если нужные детали стоят в нужном порядке, скажем, два отрезка сходятся уголком. Потом --- «или по всем местам»: следующий нейрон срабатывает, если уголок есть где угодно поблизости, неважно где именно. Первая операция делает узнавание разборчивым, вторая --- терпимым к сдвигу. Ступень за ступенью детали складываются в части, части --- в предметы, а сдвиги и масштабы прощаются.

Именно эту схему скопировали свёрточные нейросети, которые распознают картинки в телефоне. И они отвечают взаимностью: сети, обученные узнавать предметы, неплохо предсказывают, как отвечают нейроны на верхних ступенях зрения.

\section*{Учиться у видео}

Но кто учил мозг? Никто не показывал вам миллион подписанных картинок «кошка». Правдоподобный ответ --- видео. Мир течёт непрерывно, и за долю секунды предмет почти не меняется, а меняются его положение и освещение. Если связи крепнут между тем, что видно сейчас, и тем, что было мгновение назад, сеть сама свяжет разные виды одного предмета: медленное --- это предмет, быстрое --- это обстоятельства. Есть опыт, где предметы незаметно подменяли во время движения глаз, --- и нейроны начинали считать разные предметы одним. Как общее правило это пока гипотеза.

\section*{Иллюзии --- отпечатки механизмов}

Иллюзии --- не поломки, а следы того, как машина устроена. Если глаз долго смотрит на водопад, а потом на неподвижные камни, камни «ползут» вверх: пара каналов «вниз» и «вверх» разбалансирована усталостью одного из них. Знаменитое платье, которое одни видят бело-золотым, а другие сине-чёрным, делит людей по тому, какое освещение их мозг молча предполагает. Нарисованные «вращающиеся змеи» кажутся движущимися, потому что яркие участки обрабатываются чуть быстрее тусклых, и разница во времени читается как движение. А куб, который то выпирает, то проваливается, --- это две группы нейронов, соревнующиеся друг с другом; смену исхода, по одной из моделей, в основном решает шум.

Любопытно, что нейросеть, обученная предсказывать следующий кадр видео, тоже «видит», как вращаются змеи. Значит, некоторые иллюзии --- неизбежная цена за умение предсказывать.

\section*{Мозг против нейросети}

Сходство не полное. Мозгу хватает немногих примеров, он учится почти без подписей и тратит на всё двадцать ватт. Нейросеть можно обмануть незаметной человеку подделкой картинки. Правда, и людей подобные подделки немного сбивают, если показать их на мгновение. Где проходит настоящая граница между этими двумя машинами --- ещё выясняют.

\fullversion{12}
